\documentclass[11pt]{article}
\usepackage[margin=1in]{geometry}
\usepackage{amsmath,amssymb,amsthm}
\usepackage{graphicx}
\usepackage{microtype}
\usepackage[hidelinks]{hyperref}

\newtheorem{theorem}{Theorem}
\newcommand{\B}{\mathbb D}
\newcommand{\Id}{\operatorname{Id}}

\title{A Radial-Twist Counterexample for Nonlinear Abel Averages\\
\large A negative answer to the open question in arXiv:1311.6698}
\author{Automated mathematics run \texttt{fa\_banach\_001}, lane 8}
\date{July 21, 2026}

\begin{document}
\maketitle

\noindent\textbf{Status.} Candidate full counterexample, likely valid; human
review requested.

\section{Source question}

Bracci, Kozitsky, and Shoikhet \cite{BKS} ask the following. Let $D$ be a
balanced dense subset of the open unit ball of a complex Banach space, and let
$h:D\to X$ be weakly closed and $0$-dissipative. If for some
$\alpha_0>0$ the Abel average
\[
  \Phi_{\alpha_0}=(\Id-\alpha_0 h)^{-1}
\]
exists, is holomorphic, and maps the ball into itself, must the Abel averages
exist for every positive parameter?

\begin{center}
\includegraphics[width=\textwidth]{figures/open_problem_crop.png}
\end{center}

\noindent This is a full-width crop of Section 4, ``Open question,'' on page
13 of the source PDF.

\section{Idea of the counterexample}

Work in the unit disc and prescribe $\Id-h$ rather than $h$.  Inside a disc
slightly larger than radius $1/2$, let $\Id-h$ be the dilation $z\mapsto2z$.
Outside, keep its output modulus equal to $2|z|$ but gradually rotate it
through an angle of at most $\pi/3$.  The modulus $2|z|$ makes $\Id-h$
globally injective, while the absence of rotation on the preimage of the unit
disc makes its Abel average exactly $w\mapsto w/2$.

The rotation is small enough to preserve dissipativity: its cosine is at
least $1/2$.  But at parameter $1/2$, the map is the midpoint of the identity
and the twisted dilation.  Two carefully chosen circles then have the same
image radius, and their angular phases can be aligned.  This creates an exact
two-point collision over a point strictly inside the unit disc.

\section{Construction and theorem}

Let $X=\mathbb C$ with its usual norm, let $\B=\{z:|z|<1\}$, and take
$D=\B$.  Set
\[
  a=\frac{\sqrt7}{5},\qquad b=\frac35.
\]
Thus $1/2<a<b<1$.  Define a continuous function
$\varphi:[0,1)\to[0,\pi/3]$ by
\[
\varphi(r)=
\begin{cases}
0, & 0\le r\le a,\\[2mm]
\displaystyle \frac{\pi}{3}\frac{r-a}{b-a}, & a<r<b,\\[3mm]
\displaystyle \frac{\pi}{3}, & b\le r<1.
\end{cases}
\]
Finally define
\[
  h(z)=\bigl(1-2e^{i\varphi(|z|)}\bigr)z,\qquad z\in\B.
\]

\begin{theorem}
The map $h:D\to\mathbb C$ is weakly closed and $0$-dissipative.  Its Abel
average at $\alpha_0=1$ exists and is the holomorphic self-map
\[
  \Phi_1(w)=\frac{w}{2},\qquad w\in\B.
\]
Nevertheless, the Abel average at $\alpha=1/2$ does not exist: the map
$\Id-\tfrac12h$ is not injective over $\B$.  Consequently the open question
of \cite{BKS} has a negative answer.
\end{theorem}

\begin{proof}
The domain $D=\B$ is balanced and dense in $\B$.  Since $\varphi$ is
continuous, so is $h$.  In finite dimension weak and norm convergence agree,
so $h$ is weakly closed (indeed, it has a norm-closed graph).

Identify $\mathbb C^*$ with $\mathbb C$ in the usual way.  For $z\ne0$, the
unique element of $J(z)$ is the functional $z^*(w)=\overline z w$.  Therefore
\[
\operatorname{Re}\langle h(z),z^*\rangle
=\bigl(1-2\cos\varphi(|z|)\bigr)|z|^2\le0,
\]
because $0\le\varphi\le\pi/3$.  Hence the definition in \cite{BKS} holds
with $\varsigma=0$ (in fact the inequality holds throughout the disc), and
$h$ is $0$-dissipative.

At parameter $1$ we have
\[
  (\Id-h)(z)=2e^{i\varphi(|z|)}z.                 \tag{1}
\]
This map is globally injective on $\B$: equality of two outputs first gives
equality of their moduli, since the output modulus is $2|z|$; the two values
of $\varphi$ then agree, and (1) gives equality of the inputs.  In fact (1)
maps $\B$ bijectively onto the disc of radius $2$: for a point $w$ in that
disc, put $r=|w|/2$ and $z=\tfrac12e^{-i\varphi(r)}w$.

For $w\in\B$, its unique preimage under (1) has modulus $|w|/2<1/2<a$.
On that region $\varphi=0$, so the preimage is simply $w/2$.  Thus
\[
  \Phi_1=(\Id-h)^{-1}\big|_{\B}:w\longmapsto\frac w2
\]
is holomorphic and maps $\B$ into itself, exactly as required by the source
hypothesis.

It remains to show failure at parameter $1/2$.  We have
\[
  \left(\Id-\frac12h\right)(z)
  =\left(\frac12+e^{i\varphi(|z|)}\right)z.       \tag{2}
\]
Choose $\psi\in(0,\pi/2)$ by
\[
  e^{i\psi}=\frac{2+i\sqrt3}{\sqrt7},
\]
and set
\[
  z_1=\frac{\sqrt7}{5},\qquad
  z_2=\frac35e^{-i\psi}
      =\frac{3(2-i\sqrt3)}{5\sqrt7}.
\]
Their moduli are $a$ and $b$, respectively, so
$\varphi(|z_1|)=0$ and $\varphi(|z_2|)=\pi/3$.  Formula (2) gives
\[
\left(\Id-\frac12h\right)(z_1)
=\frac32\frac{\sqrt7}{5}
=\frac{3\sqrt7}{10}.
\]
Also
\[
  \frac12+e^{i\pi/3}
  =\frac{2+i\sqrt3}{2}
  =\frac{\sqrt7}{2}e^{i\psi},
\]
and hence
\[
\left(\Id-\frac12h\right)(z_2)
=\frac{\sqrt7}{2}e^{i\psi}\frac35e^{-i\psi}
=\frac{3\sqrt7}{10}.
\]
The points $z_1$ and $z_2$ are distinct because their moduli differ, whereas
their common image belongs to $\B$ since
\[
  \left|\frac{3\sqrt7}{10}\right|^2=\frac{63}{100}<1.
\]
Thus $(\Id-\tfrac12h)^{-1}$ is not a function on $\B$, so the Abel average
at $\alpha=1/2$ does not exist.  This proves the negative answer.
\end{proof}

\section{Checks, scope, and relation to the partial theorem}

The construction satisfies a stronger property than the source hypothesis:
$\Id-h$ is globally injective on $D$ and even maps $D$ bijectively onto the
disc of radius $2$.  Thus the counterexample is insensitive to whether the
inverse notation is read globally or only over the unit ball.

The earlier upward-propagation result for this target remains consistent with
the example: existence at $\alpha_0=1$ forces all averages for
$\alpha\ge1$, while the present failure occurs at $\alpha=1/2$.  The example
therefore shows that the unresolved downward direction is a genuine
obstruction and settles the full source question negatively.

The optional script \texttt{code/check\_counterexample.py} samples the
dissipativity inequality and evaluates the collision numerically.  It is not
part of the proof; all calculations above are exact.

\section{Novelty check and review recommendation}

A bounded search covered the exact printed question, arXiv id 1311.6698, the
source title and authors, and combinations of ``weakly closed,''
``0-dissipative,'' ``Abel average,'' ``counterexample,'' and ``negative
answer.''  No later answer or matching construction was found.  The 2019
monograph \cite{ERS} reproduces the question as open.  Novelty confidence is
therefore moderate-high within the searched sources, but exhaustive priority
is not claimed.

Human review should focus on the support-functional calculation, global
injectivity of (1), and the two exact evaluations of (2).  No unproved lemma
or external theorem is used.

\begin{thebibliography}{9}
\bibitem{BKS}
F. Bracci, Y. Kozitsky, and D. Shoikhet,
\emph{Abel averages and holomorphically pseudo-contractive maps in Banach
spaces}, J. Math. Anal. Appl. 423 (2015), no. 2, 1580--1593;
arXiv:1311.6698. DOI: 10.1016/j.jmaa.2014.10.079.

\bibitem{ERS}
M. Elin, S. Reich, and D. Shoikhet,
\emph{Numerical Range of Holomorphic Mappings and Applications},
Birkh\"auser, Cham, 2019. DOI: 10.1007/978-3-030-05020-7.
\end{thebibliography}

\end{document}
